% Auto-generated by citations/build_ledger_pdf.py — do not hand-edit.
\documentclass[11pt]{article}
\usepackage{lmodern}
\usepackage[margin=1in]{geometry}
\usepackage[T1]{fontenc}
\usepackage[utf8]{inputenc}
\usepackage{microtype}
\usepackage{xurl}
\usepackage{xcolor}
\usepackage{enumitem}
\usepackage[colorlinks=true,urlcolor=blue!50!black,linkcolor=black]{hyperref}
\hypersetup{
  pdftitle={esperanto-research: Grounded Citation Ledger},
  pdfsubject={Citation ledger for the OPUS23 | ESPERANTO grounded review (51 registered sources)},
  pdfauthor={Hermes Agent + Walker (grounded-citations workflow)},
  pdfkeywords={grounded-citations, citation ledger, OPUS23, ESPERANTO, Datapunk, evidence quotes}
}
\setlength{\parindent}{0pt}
\setlength{\parskip}{6pt}
\begin{document}

\begin{center}
{\LARGE \textbf{Grounded Citation Ledger}}\\[6pt]
{\large esperanto-research --- a citation-verified review of OPUS23 \textbar{} ESPERANTO}\\[4pt]
{\small Generated 2026-10-30 from \texttt{citations/ledger.json} (grounded-citations workflow)}
\end{center}

\section*{Method}
Every citation in this project's reports resolves to an entry in the mechanical ledger (file
\path{citations/ledger.json}); ids are issued only at retrieval, never from memory.
This document renders the full ledger: \textbf{51 registered sources},
of which the [1, 2, 4, 5] platform sources carry \textbf{16 verbatim evidence quotes}
--- each machine-checked (``\texttt{sources.py quote -{}-from}'') to appear verbatim,
whitespace- and markup-tolerant, in the fetched page text preserved under \path{extracts/}.
All pages were accessed 2026-10-30.
Quotation text below is reproduced exactly as attached (raw page markup such as \texttt{**}
may be retained inside quotes). Regenerate with \path{citations/build_ledger_pdf.py}.

\section*{Platform sources --- evidence-carrying}
\noindent\textbf{[1]} OPUS23 | ESPERANTO \textperiodcentered{} Datapunk Bioinformatics (platform landing page)\\
\indent\url{https://www.datapunk.net} \emph{(accessed 2026-10-30)}
\\ \textbf{Verbatim evidence:}
\quad---\,\emph{``A 22,800-taxon knowledge base behind every stool report.''}
\quad---\,\emph{``390 explained markers, 3,200+ metabolites''}
\quad---\,\emph{``Drafted by a dedicated large language model in the background; every claim carries a citation you can click.''}
\quad---\,\emph{``Opie is our literature agent.''}
\quad---\,\emph{``Pay as you go, with tokens.''}
\quad---\,\emph{``It does not diagnose, treat or prescribe''}
\quad---\,\emph{``ABO, secretor (FUT2) and Lewis status''}
\quad---\,\emph{``389** multi-SNP algorithms''}
\quad---\,\emph{``4,500+** curated SNPs''}
\vspace{4pt}

\noindent\textbf{[2]} Datapunk Opus23 Pro (legacy splash)\\
\indent\url{https://www.datapunk.net/opus23} \emph{(accessed 2026-10-30)}
\\ \textbf{Verbatim evidence:}
\quad---\,\emph{``Over 200 clinicians worldwide use Opus23''}
\quad---\,\emph{``Each license is for one client''}
\quad---\,\emph{``report over 300 potential adverse reactions''}
\vspace{4pt}

\noindent\textbf{[3]} Datapunk Open Source Apps\\
\indent\url{https://datapunk.net/tlfd} \emph{(accessed 2026-10-30)}
\vspace{4pt}

\noindent\textbf{[4]} Opus23 Pro Help Docs --- Welcome\\
\indent\url{https://www.datapunk.net/opus23/docs/index.pl?WELCOME} \emph{(accessed 2026-10-30)}
\\ \textbf{Verbatim evidence:}
\quad---\,\emph{``runs over 300 algorithms on multi-snp clusters''}
\quad---\,\emph{``exceeds HIPAA standards''}
\quad---\,\emph{``patients at the University of Bridgeport, CT''}
\vspace{4pt}

\noindent\textbf{[5]} Opus23 Developer Blog --- About Opus 23\\
\indent\url{https://datapunk.net/opus23blog/about_opus23} \emph{(accessed 2026-10-30)}
\\ \textbf{Verbatim evidence:}
\quad---\,\emph{``Center of Excellence in Generative Medicine (COEGM) to analyze and interpret raw genomic data''}
\vspace{4pt}

\section*{Aggregated dossier sources --- competitor verification (03)}
\noindent\textbf{[6]} Jona Clinical Gut Microbiome Analysis Advanced\\
\indent\url{https://jona.health/products/jona-clinical-gut-microbiome-analysis-advanced} \emph{(accessed 2026-10-30)}

\noindent\textbf{[7]} DNA Testing \& Precision Medicine Platform for Medical Doctors\\
\indent\url{https://promicsedge.com/medical-doctors} \emph{(accessed 2026-10-30)}

\noindent\textbf{[8]} Lite - OmicsEdge\\
\indent\url{https://promicsedge.com/lite} \emph{(accessed 2026-10-30)}

\noindent\textbf{[9]} Store - Tiny Health\\
\indent\url{https://www.tinyhealth.com/store} \emph{(accessed 2026-10-30)}

\noindent\textbf{[10]} Core Panels | Fagron Genomics US\\
\indent\url{https://gxsciences.com/nutrigenomics} \emph{(accessed 2026-10-30)}

\noindent\textbf{[11]} Pro7 Panel | Fagron Genomics US\\
\indent\url{https://gxsciences.com/nutrigenomics/pro7-panel} \emph{(accessed 2026-10-30)}

\noindent\textbf{[12]} Best AI Scribe for Functional Medicine (2026) | DeepCura Resources\\
\indent\url{https://www.deepcura.com/resources/best-ai-scribe-for-functional-medicine} \emph{(accessed 2026-10-30)}

\noindent\textbf{[13]} DeepCura --- Beyond the EHR: The Electronic Health Workforce\\
\indent\url{https://www.deepcura.com} \emph{(accessed 2026-10-30)}

\noindent\textbf{[14]} Microbiome Analysis Software - Metabolon\\
\indent\url{https://www.metabolon.com/bioinformatics-software/microbiome-software} \emph{(accessed 2026-10-30)}

\noindent\textbf{[15]} 3MCor: An Integrative Web Server for Metabolome-Microbiome-Metadata Correlation Analysis (preprint)\\
\indent\url{https://pdfs.semanticscholar.org/e7bc/6cd38db249ae0f9448dc29195d3c8ef52cb5.pdf} \emph{(accessed 2026-10-30)}

\noindent\textbf{[16]} AmiAge Clock\\
\indent\url{https://3mcor.cn} \emph{(accessed 2026-10-30)}

\noindent\textbf{[18]} MMIP --- Microbiome Metabolome Integration Platform\\
\indent\url{https://bioinfo.jisiasr.org/mmip/index.html} \emph{(accessed 2026-10-30)}

\noindent\textbf{[19]} FDX Method - Functional DX\\
\indent\url{https://functionaldx.com/fdx-method} \emph{(accessed 2026-10-30)}

\noindent\textbf{[20]} Home Functional DX\\
\indent\url{https://functionaldx.com} \emph{(accessed 2026-10-30)}

\section*{Aggregated dossier sources --- science-flags literature check (05)}
\noindent\textbf{[17]} Microbiome Metabolome Integration Platform (MMIP) --- PubMed PMID 37771003\\
\indent\url{https://pubmed.ncbi.nlm.nih.gov/37771003} \emph{(accessed 2026-10-30)}

\noindent\textbf{[21]} Identification of human zonulin, a physiological modulator of tight junctions, as prehaptoglobin-2 (PNAS 2009)\\
\indent\url{https://pubmed.ncbi.nlm.nih.gov/19805376} \emph{(accessed 2026-10-30)}

\noindent\textbf{[22]} Zonulin, regulation of tight junctions, and autoimmune diseases (Ann NY Acad Sci 2012)\\
\indent\url{https://pubmed.ncbi.nlm.nih.gov/22731712} \emph{(accessed 2026-10-30)}

\noindent\textbf{[23]} Celiac disease (Curr Opin Gastroenterol 2009)\\
\indent\url{https://pubmed.ncbi.nlm.nih.gov/19122520} \emph{(accessed 2026-10-30)}

\noindent\textbf{[24]} A purified membrane protein from Akkermansia muciniphila or the pasteurized bacterium improves metabolism in obese and diabetic mice (Nat Med 2017)\\
\indent\url{https://pubmed.ncbi.nlm.nih.gov/27892954} \emph{(accessed 2026-10-30)}

\noindent\textbf{[25]} Commensal Akkermansia muciniphila exacerbates gut inflammation in Salmonella Typhimurium-infected gnotobiotic mice (PLoS ONE 2013)\\
\indent\url{https://pubmed.ncbi.nlm.nih.gov/24040367} \emph{(accessed 2026-10-30)}

\noindent\textbf{[26]} Faecalibacterium prausnitzii is an anti-inflammatory commensal bacterium identified by gut microbiota analysis of Crohn disease patients (PNAS 2008)\\
\indent\url{https://pubmed.ncbi.nlm.nih.gov/18936492} \emph{(accessed 2026-10-30)}

\noindent\textbf{[27]} Bifidobacteria and Butyrate-Producing Colon Bacteria: Importance and Strategies for Their Stimulation in the Human Gut (Front Microbiol 2016)\\
\indent\url{https://pubmed.ncbi.nlm.nih.gov/27446020} \emph{(accessed 2026-10-30)}

\noindent\textbf{[28]} Widely Used Commercial ELISA Does Not Detect Precursor of Haptoglobin2, but Recognizes Properdin as a Potential Second Member of the Zonulin Family (Front Endocrinol 2018)\\
\indent\url{https://pubmed.ncbi.nlm.nih.gov/29459849} \emph{(accessed 2026-10-30)}

\noindent\textbf{[29]} A fast, cheap, and easy protocol for celiac disease HLA haplotype typing using buccal swabs (Methods Cell Biol 2023)\\
\indent\url{https://pubmed.ncbi.nlm.nih.gov/37625876} \emph{(accessed 2026-10-30)}

\noindent\textbf{[30]} Effective detection of human leukocyte antigen risk alleles in celiac disease using tag single nucleotide polymorphisms (PLoS ONE 2008)\\
\indent\url{https://pubmed.ncbi.nlm.nih.gov/18509540} \emph{(accessed 2026-10-30)}

\noindent\textbf{[31]} Cost-effective HLA typing with tagging SNPs predicts celiac disease risk haplotypes in the Finnish, Hungarian, and Italian populations (Immunogenetics 2009)\\
\indent\url{https://pubmed.ncbi.nlm.nih.gov/19255754} \emph{(accessed 2026-10-30)}

\noindent\textbf{[32]} ESPGHAN guidelines for the diagnosis of coeliac disease (J Pediatr Gastroenterol Nutr 2012)\\
\indent\url{https://pubmed.ncbi.nlm.nih.gov/22197856} \emph{(accessed 2026-10-30)}

\noindent\textbf{[33]} ESPGHAN Guidelines for Diagnosing Coeliac Disease 2020 (J Pediatr Gastroenterol Nutr)\\
\indent\url{https://pubmed.ncbi.nlm.nih.gov/31568151} \emph{(accessed 2026-10-30)}

\noindent\textbf{[34]} The GenoChip: a new tool for genetic anthropology (Genome Biol Evol 2013)\\
\indent\url{https://pubmed.ncbi.nlm.nih.gov/23666864} \emph{(accessed 2026-10-30)}

\noindent\textbf{[35]} Mitochondrial DNA Haplogroup K Is Protective Against Autism Spectrum Disorder Risk (JAACAP 2024)\\
\indent\url{https://pubmed.ncbi.nlm.nih.gov/38072244} \emph{(accessed 2026-10-30)}

\noindent\textbf{[36]} Examining the role of common and rare mitochondrial variants in schizophrenia (PLoS ONE 2018)\\
\indent\url{https://pubmed.ncbi.nlm.nih.gov/29370225} \emph{(accessed 2026-10-30)}

\noindent\textbf{[37]} The association between common vitamin D receptor gene variations and osteoporosis: a participant-level meta-analysis (Ann Intern Med 2006)\\
\indent\url{https://pubmed.ncbi.nlm.nih.gov/16908916} \emph{(accessed 2026-10-30)}

\noindent\textbf{[38]} BsmI, TaqI, ApaI and FokI polymorphisms in the vitamin D receptor (VDR) gene and risk of fracture in Caucasians: a meta-analysis (Bone 2011)\\
\indent\url{https://pubmed.ncbi.nlm.nih.gov/20601302} \emph{(accessed 2026-10-30)}

\noindent\textbf{[39]} Do Vitamin D receptor gene polymorphisms affect bone mass density in men? (Ageing Res Rev 2022)\\
\indent\url{https://pubmed.ncbi.nlm.nih.gov/35063697} \emph{(accessed 2026-10-30)}

\noindent\textbf{[40]} Common genetic determinants of vitamin D insufficiency: a genome-wide association study (Lancet 2010)\\
\indent\url{https://pubmed.ncbi.nlm.nih.gov/20541252} \emph{(accessed 2026-10-30)}

\noindent\textbf{[41]} A second common mutation in the methylenetetrahydrofolate reductase gene: an additional risk factor for neural-tube defects? (Am J Hum Genet 1998)\\
\indent\url{https://pubmed.ncbi.nlm.nih.gov/9545395} \emph{(accessed 2026-10-30)}

\noindent\textbf{[42]} MTHFR 677C-->T polymorphism and risk of coronary heart disease: a meta-analysis (JAMA 2002)\\
\indent\url{https://pubmed.ncbi.nlm.nih.gov/12387655} \emph{(accessed 2026-10-30)}

\noindent\textbf{[43]} ACMG Practice Guideline: lack of evidence for MTHFR polymorphism testing (Genet Med 2013)\\
\indent\url{https://pubmed.ncbi.nlm.nih.gov/23288205} \emph{(accessed 2026-10-30)}

\noindent\textbf{[44]} Proposed guidelines to evaluate scientific validity and evidence for genotype-based dietary advice (Genes Nutr 2017)\\
\indent\url{https://pubmed.ncbi.nlm.nih.gov/29270237} \emph{(accessed 2026-10-30)}

\noindent\textbf{[45]} APOA2, dietary fat, and body mass index: replication of a gene-diet interaction in 3 independent populations (Arch Intern Med 2009)\\
\indent\url{https://pubmed.ncbi.nlm.nih.gov/19901143} \emph{(accessed 2026-10-30)}

\noindent\textbf{[46]} Apolipoprotein A2 polymorphism interacts with intakes of dairy foods to influence body weight in 2 U.S. populations\\
\indent\url{https://pubmed.ncbi.nlm.nih.gov/24108135} \emph{(accessed 2026-10-30)}

\noindent\textbf{[47]} Adiposity-related evidence for APOA2, FTO, and MC4R: a PRISMA-ScR scoping review of recent human studies (2019-2025)\\
\indent\url{https://pubmed.ncbi.nlm.nih.gov/42950483} \emph{(accessed 2026-10-30)}

\noindent\textbf{[48]} Genes involved in the regulation of intestinal permeability and their role in ulcerative colitis (J Dig Dis 2016)\\
\indent\url{https://pubmed.ncbi.nlm.nih.gov/26512799} \emph{(accessed 2026-10-30)}

\noindent\textbf{[49]} Polymorphisms in E-cadherin (CDH1) result in a mis-localised cytoplasmic protein that is associated with Crohn's disease (Gut 2009)\\
\indent\url{https://pubmed.ncbi.nlm.nih.gov/19398441} \emph{(accessed 2026-10-30)}

\noindent\textbf{[50]} Polygenic Risk Scores for Cardiovascular Disease: A Scientific Statement From the American Heart Association (Circulation 2022)\\
\indent\url{https://pubmed.ncbi.nlm.nih.gov/35862132} \emph{(accessed 2026-10-30)}

\noindent\textbf{[51]} Clinical use of current polygenic risk scores may exacerbate health disparities (Nat Genet 2019)\\
\indent\url{https://pubmed.ncbi.nlm.nih.gov/30926966} \emph{(accessed 2026-10-30)}

\end{document}
